\documentclass[12pt]{article}
\usepackage[utf8]{inputenc}
\usepackage[T1]{fontenc}
\usepackage{amsmath, amssymb}
\usepackage{geometry}
\usepackage{xcolor}
\usepackage{lipsum}
\usepackage{titlesec}
\usepackage{fancyhdr}
\usepackage{tcolorbox}
\usepackage{graphicx}
\usepackage{float}
\usepackage{tikz}
\usetikzlibrary{shapes.geometric, arrows.meta}
\usetikzlibrary{arrows.meta, decorations.markings}

\geometry{a4paper, margin=2.5cm}
\pagestyle{fancy}
\fancyhf{}
\rhead{Problem Set I}
\lhead{Analysis IV}
\cfoot{\thepage}

% Fancy titles
\titleformat{\section}{\normalfont\Large\bfseries}{Exercise \thesection:}{1em}{}
\titleformat{\subsection}{\normalfont\bfseries}{Correction:}{1em}{}

% Custom box for correction
\newtcolorbox{correctionbox}{
  colback=gray!5,
  colframe=black,
  fonttitle=\bfseries,
  title=Correction,
  before skip=10pt,
  after skip=10pt
}

% Fix headheight warning
\setlength{\headheight}{14.49998pt}

\begin{document}

\begin{center}
\Large\textbf{Ibn Tofail University} \\[1em]
\large\textit{Analysis IV} \\[0.5em]
\large\textit{Problem Set I} \\[0.5em]
\end{center}

\vspace{1cm}

% ----------- EXERCISE 1 -------------
\section{}
Study the convergence of the following series:

\begin{enumerate}
    \item \( \displaystyle \sum_{n=1}^{+\infty} \frac{1}{n} \)
    
    \item \( \displaystyle \sum_{n=2}^{+\infty} \frac{n^2 + 1}{n^2} \)
    
    \item \( \displaystyle \sum_{n=1}^{+\infty} \frac{2}{\sqrt{n}} \)
    
    \item \( \displaystyle \sum_{n=1}^{+\infty} \frac{(2n + 1)^4}{(7n^2 + 1)^3} \)
    
    \item \( \displaystyle \sum_{n=1}^{+\infty} \left( ne^{\frac{1}{n}} - n \right) \)
    
    \item \( \displaystyle \sum_{n=0}^{+\infty} \ln\left(1 + e^{-n}\right) \)
    
    \item \( \displaystyle \sum_{n=1}^{+\infty} \left( \frac{n}{n + 1} \right)^{n^2} \)
\end{enumerate}

\begin{correctionbox}
Solution here.
\end{correctionbox}

% ----------- EXERCISE 2 -------------
\section{}
Determine the nature of the numerical series with the following general terms:

\begin{enumerate}
    \item \( u_n = \frac{1}{n \cos^2 n} \)
    
    \item \( v_n = \frac{1}{[\ln(n)]^n} \)
    
    \item \( w_n = \frac{\tan^n\left(\frac{\pi}{7}\right)}{3^{n+2}} \)
    
    \item \( x_n = \frac{1}{n \ln(n)} \)
    
    \item \( y_n = \frac{1}{n(\ln n)^2} \)
    
    \item \( z_n = \frac{1}{(\ln 2)^2 + (\ln 3)^2 + \ldots + (\ln n)^2} \)
\end{enumerate}

\begin{correctionbox}
Solution here.
\end{correctionbox}

% ----------- EXERCISE 3 -------------
\section{}
Show that the series with general term
\[
w_n = \frac{(-1)^n}{\ln(1 + \sqrt{n})}
\]
is convergent but not absolutely convergent (called semi-convergent).

\begin{correctionbox}
Solution here.
\end{correctionbox}

% ----------- EXERCISE 4 -------------
\section{}
\begin{enumerate}
    \item Show that the following two series are absolutely convergent:
    \[
    \sum_{n=0}^{+\infty} \frac{1}{n!} \quad \text{and} \quad \sum_{n=0}^{+\infty} \frac{(-1)^n}{2^n}
    \]
    
    \item Deduce that the series:
    \[
    \sum_{n=0}^{+\infty} \left[ \sum_{k=0}^{n} \frac{(-1)^{n-k}}{k! \cdot 2^{n-k}} \right]
    \]
    is absolutely convergent and determine its value.
\end{enumerate}

\begin{correctionbox}
Solution here.
\end{correctionbox}

% ----------- EXERCISE 5 -------------
\section{}
Study the nature of the following series:
\[
\sum_{n>2} \ln\left(1 + \frac{(-1)^n}{n}\right), \quad \sum_{n>1} \sin\left(\frac{(-1)^n}{n}\right)
\]

\begin{correctionbox}
Solution here.
\end{correctionbox}

% ----------- EXERCISE 6 -------------
\section{}
Show that the series with general terms
\[
u_n := \frac{(-1)^n}{\sqrt{n}} \quad \text{and} \quad v_n := \frac{(-1)^n}{\sqrt{n + (-1)^n}}
\]
are not of the same nature, although \( u_n \sim v_n \).

\begin{correctionbox}
Solution here.
\end{correctionbox}

\end{document}